\documentclass[12pt]{article}
\usepackage[T1]{fontenc}
\usepackage[utf8]{inputenc}
\usepackage{lmodern}
\usepackage{textcomp}
\usepackage{enumitem}
\usepackage{geometry}
\geometry{margin=1in}
\begin{document}
\begin{center}
Answer Key - Psychology - Undergraduate Level Assessment 13 \
\small (Answers are ordered exactly as in the question paper.)
\end{center}

\section*{Multiple Choice}
\begin{enumerate}[label=\arabic*.]
\item \textbf{Answer:} D
\\ \textit{Options:} A) occipital lobe; B) temporal lobe; C) parietal lobe; D) cerebral cortex; E) limbic system
\\ \textit{Note:} The cerebral cortex is the least developed area of the brain at birth because it undergoes significant growth and maturation during infancy and early childhood, while other brain structures, such as the brainstem and cerebellum, are more fully developed and functioning at birth.
\item \textbf{Answer:} B
\\ \textit{Options:} A) 0; B) 0.5; C) 0.8; D) 0.3; E) 1.2
\\ \textit{Note:} A difficulty level (p) of 0.5 maximizes differentiation between high- and low-performing groups because it indicates that approximately half of the examinees answer the item correctly, providing the greatest potential for distinguishing between varying levels of knowledge or ability.
\item \textbf{Answer:} E
\\ \textit{Options:} A) Universal, selective, and preventive; B) Primary, secondary, and indicative; C) Primary, secondary, and quaternary prevention; D) General, selective, and indicative; E) Universal, selective, and indicative
\\ \textit{Note:} The correct answer is based on the understanding that primary prevention encompasses three distinct approaches: universal prevention targets the entire population, selective prevention focuses on at-risk groups, and indicative prevention addresses individuals who are already showing early signs of a problem, thereby forming a comprehensive strategy to prevent issues before they arise.
\item \textbf{Answer:} A
\\ \textit{Options:} A) phonemes from the child's native language only.; B) Phonemes from the language most frequently spoken around the child.; C) Phonetic sounds specific to the child's geographical location.; D) Phonemes from the parents' native language only.; E) morphemes that the child has heard most frequently.
\\ \textit{Note:} The correct answer is supported by the fact that during the babbling stage, infants begin to produce sounds that are primarily influenced by the phonetic patterns of their native language, gradually narrowing their vocalizations to the specific phonemes they are exposed to in their linguistic environment.
\item \textbf{Answer:} A
\\ \textit{Options:} A) Have no significant effect on the performance of tasks; B) Facilitate the performance of both easy and difficult tasks; C) hinder the performance of all tasks; D) Hinder the performance of difficult tasks; E) hinder the performance of easy tasks
\\ \textit{Note:} The correct answer reflects the concept of social facilitation and inhibition, indicating that while the presence of others can enhance performance on simple tasks, it does not significantly affect performance on complex or unfamiliar tasks, leading to an overall conclusion that, in many scenarios, there is no significant effect on performance.
\end{enumerate}

\section*{True / False}
\begin{enumerate}[label=\arabic*.]
\setcounter{enumi}{5}
\item \textbf{Answer:} False
\\ \textit{Options:} A) True; B) False
\\ \textit{Note:} The Solomon four-group design is primarily intended to control for threats to internal validity, particularly testing effects, rather than to increase sample size.
\item \textbf{Answer:} True
\\ \textit{Options:} A) True; B) False
\\ \textit{Note:} Carlos could differentiate middle C on the violin from middle C on the piano because each instrument produces a unique timbre, which encompasses the specific qualities of sound that allow listeners to identify different sources even when the pitch is the same.
\item \textbf{Answer:} False
\\ \textit{Options:} A) True; B) False
\\ \textit{Note:} In Bronfenbrenner's ecological model, the nanosystem is not a recognized term; instead, the microsystem refers to the direct interactions between individuals, such as family and school, while the nanosystem is not a level within this framework.
\item \textbf{Answer:} True
\\ \textit{Options:} A) True; B) False
\\ \textit{Note:} Individuals with antisocial personalities often prioritize immediate gratification and tangible rewards, making money a compelling motivator for their learning and behavior.
\item \textbf{Answer:} False
\\ \textit{Options:} A) True; B) False
\\ \textit{Note:} Strict punishments can create fear and resentment, potentially leading to decreased prosocial behavior, while positive reinforcement and fostering a supportive community are more effective in promoting helping behaviors.
\end{enumerate}

\section*{Long Form Response}
\begin{enumerate}[label=\arabic*.]
\setcounter{enumi}{10}
\item \textbf{Answer:} The medulla, located at the base of the brainstem, plays a crucial role in the central nervous system by regulating vital bodily functions. It is responsible for autonomic processes such as heart rate, breathing, and blood pressure, ensuring that these essential functions occur without conscious effort. The medulla also serves as a pathway for communication between the brain and spinal cord, facilitating reflex actions and maintaining homeostasis. Its significance is underscored by the fact that damage to the medulla can lead to life-threatening conditions, highlighting its integral role in sustaining life. Overall, the medulla is essential for maintaining the physiological balance necessary for survival.
\\ \textit{Note:} correct because it accurately identifies the medulla's location at the brainstem, its critical functions in regulating autonomic processes essential for survival, and its role in facilitating communication within the central nervous system, thus highlighting its significance in maintaining homeostasis.
\item \textbf{Answer:} Self-monitoring refers to the degree to which individuals regulate their behavior in social contexts based on external cues and the reactions of others. Individuals low in self-monitoring tend to rely more heavily on life scripts-internalized narratives that guide their behavior in various situations-rather than adapting their actions to fit social expectations. These life scripts are shaped by personal experiences and cultural contexts, leading low self-monitors to act consistently with their established beliefs and values, regardless of the social setting. Consequently, their behavior may appear more authentic but can also lead to social misalignments or misunderstandings, as they may not adjust to the nuances of varying social interactions. This reliance on life scripts highlights the contrast between self-monitoring strategies and personal authenticity in social behavior.
\\ \textit{Note:} correct because it accurately describes self-monitoring as the regulation of behavior based on social cues and explains how individuals low in self-monitoring rely on life scripts, which are shaped by personal and cultural experiences, to guide their behavior consistently rather than adapting to social expectations.
\end{enumerate}

\vfill
\noindent\textit{End of Answer Key}
\end{document}